\documentclass[12pt]{article}
\usepackage[a4paper,margin=1in]{geometry}
\usepackage[utf8]{inputenc}
\usepackage[T1]{fontenc}
\usepackage{amsmath,amssymb}
\usepackage{textcomp}
\usepackage{parskip}
\begin{document}
\begin{center}
{\Large \textbf{2025 STEP III}}
\end{center}

\subsection*{Page 1}
\begin{flushleft}
\bigskip\textbf{Section A: Pure Mathematics}\par
\bigskip\textbf{Question 1.} You need not consider the convergence of the improper integrals in this question.\par
\[For p, q > 0, define\]
\[b(p, q) = xp 1(1 x)q 1 dx.\]
- -\par
-\par
\textbf{(i)} Show that b(p, q) = b(q, p).\par
\textbf{(ii)} Show that b(p + 1, q) = b(p, q) b(p, q + 1) and hence that b(p + 1, p) = 1 b(p, p).\par
- 2\par
\textbf{(iii)} Show that\par
\[1 \pi\]
\[b(p, q) = 2 2 (sin θ)2p 1(cos θ)2q 1 dθ.\]
- -\par
\[Hence show that b(p, p) = b p, 1 .\]
22p 1 2\par
-\par
( )\par
\textbf{(iv)} Show that\par
tp 1\par
\infty -\par
\[b(p, q) = dt.\]
(1 + t)p+q\par
\textbf{(v)} Evaluate\par
\infty t 2\par
dt.\par
(1 + t)6\par
\end{flushleft}
\subsection*{Page 2}
\begin{flushleft}
\bigskip\textbf{Question 2.} Let f(x) = 7 2 x .\par
- | |\par
\[A sequence u , u , u , . . . is defined by u = a and u = f(u ) for n > 0.\]
0 1 2 0 n n 1\par
-\par
\textbf{(i)} (a) Sketch, on the same axes, the graphs with equations y = f(x) and y = f(f(x)).\par
\[(b) Find all solutions of the equation f(f(x)) = x.\]
(c) Find the values of a for which the sequence u , u , u , . . . has period 2.\par
0 1 2\par
\[(d) Show that, if a = 28 , then the sequence u , u , u , . . . has period 2, but neither\]
5 2 3 4\par
u or u is equal to either of u or u .\par
0 1 2 3\par
\textbf{(ii)} (a) Sketch, on the same axes, the graphs with equations y = f(x) and y = f(f(f(x))).\par
\[(b) Consider the sequence u , u , u , . . . in the cases a = 1 and a = 7 . Hence find\]
0 1 2 - 9\par
\[all the solutions of the equation f(f(f(x))) = x.\]
(c) Find a value of a such that the sequence u , u , u , . . . has period 3, but where\par
3 4 5\par
none of u , u or u is equal to any of u , u or u .\par
0 1 2 3 4 5\par
\end{flushleft}
\subsection*{Page 3}
\begin{flushleft}
\bigskip\textbf{Question 3.} Let f(x) be defined and positive for x > 0.\par
\[Let a and b be real numbers with 0 < a < b and define the points A = (a, f(a)) and\]
\[B = (b, f(b)).\]
-\par
\[Let X = (m, 0) be the point of intersection of line AB with the x-axis.\]
\textbf{(i)} Find an expression for m in terms of a, b, f(a) and f(b).\par
\textbf{(ii)} Show that, if f(x) = \sqrtx, then m = \sqrtab.\par
an+1 + bn+1\par
\[Find, in terms of n, a function f(x) such that m = .\]
an + bn\par
\textbf{(iii)} Let g (x) and g (x) be defined and positive for x > 0. Let m = M when f(x) = g (x)\par
1 2 1 1\par
\[and let m = M when f(x) = g (x).\]
2 2\par
g (x)\par
\[Show that if is a decreasing function then M > M .\]
1 2\par
g (x)\par
Hence show that\par
a + b 2ab\par
\[> \sqrt\{ab\} > .\]
2 a + b\par
\textbf{(iv)} Let p and c be chosen so that the curve y = p(c x)3 passes through both A and B.\par
-\par
Show that\par
c a f(a) 3\par
\[- =\]
b c f(b)\par
-  (\par
and hence determine c in terms of a, b, f(a) and f(b).\par
\[Show that if f is a decreasing function, then c < m.\]
\end{flushleft}
\subsection*{Page 4}
\begin{flushleft}
4 (i) x and y are defined in terms of x and y by the equation\par
2 2 1 1\par
x 1 1 x\par
\[2 \sqrt\{2\} - \sqrt\{2\} 1\]
\[= .\]
\[y 2 (  \sqrt\{1\} 2 \sqrt\{1\} 2 (  y 1 (\]
G is the graph with equation\par
\[x^2 y^2\]
\[+ = 1\]
9 4\par
and G is the graph with equation\par
2 2\par
x y x y\par
+ +\par
\[\sqrt\{2\} \sqrt\{2\} - \sqrt\{2\} \sqrt\{2\}\]
\[) ) + ) ) = 1.\]
9 4\par
Show that, if (x , y ) is a point on G , then (x , y ) is a point on G .\par
1 1 1 2 2 2\par
Show that G is an anti-clockwise rotation of G through 45 about the origin.\par
2 1 \textasciicircum{}\circ\par
\textbf{(ii)} (a) The matrix\par
0.6 0.8\par
-\par
0.8 0.6\par
) )\par
represents a reflection. Find the line of invariant points of this matrix.\par
(b) Sketch, on the same axes, the graphs with equations\par
\[y = 2x and 0.8x + 0.6y = 2 0.6x+0.8y.\]
-\par
\textbf{(iii)} Sketch, on the same axes, for 0  x  2\pi, the graphs with equations\par
\[y = sin x and y = sin(x 2y).\]
-\par
You should determine the exact co-ordinates of the points on the graph with equation\par
\[y = sin(x 2y) where the tangent is horizontal and those where it is vertical.\]
-\par
\end{flushleft}
\subsection*{Page 5}
\begin{flushleft}
\bigskip\textbf{Question 5.} Three points, A, B and C, lie in a horizontal plane, but are not collinear. The point O lies\par
above the plane.\par
\[Let -O\toA = a, -O-\toB = b and -O-\toC = c.\]
P is a point with -O-\toP = αa + βb + γc, where α, β and γ are all positive and α + β + γ < 1.\par
\[Let k = 1 (α + β + γ).\]
-\par
\textbf{(i)} The point L is on OA, the point X is on BC and LX passes through P .\par
α\par
\[Determine -O-\toX in terms of β, γ, b and c and show that -O\toL = a.\]
k + α\par
\textbf{(ii)} Let M and Y be the unique pair of points on OB and CA respectively such that\par
M Y passes through P , and let N and Z be the unique pair of points on OC and AB\par
respectively such that N Z passes through P .\par
Show that the plane LM N is also horizontal if and only if OP intersects plane ABC\par
at the point G, where -O-\toG 1 (a + b + c). Where do points X, Y and Z lie in this\par
\equiv 3\par
case?\par
\textbf{(iii)} State what the condition α + β + γ < 1 tells you about the position of P relative to\par
the tetrahedron OABC.\par
\bigskip\textbf{Question 6.} (i) Let a, b and c be three non-zero complex numbers with the properties a + b + c = 0\par
\[and a^2 + b^2 + c^2 = 0.\]
Show that a, b and c cannot all be real.\par
Show further that a, b and c all have the same modulus.\par
\textbf{(ii)} Show that it is not possible to find three non-zero complex numbers a, b and c with\par
\[the properties a + b + c = 0 and a^3 + b^3 + c^3 = 0.\]
\textbf{(iii)} Show that if any four non-zero complex numbers a, b, c and d have the properties\par
\[a + b + c + d = 0 and a^3 + b^3 + c^3 + d^3 = 0, then at least two of them must have the\]
same modulus.\par
\textbf{(iv)} Show, by taking c = 1, d = 2 and e = 3 that it is possible to find five real numbers\par
-\par
\[a, b, c, d and e with distinct magnitudes and with the properties a + b + c + d + e = 0\]
\[and a^3 + b^3 + c^3 + d^3 + e^3 = 0.\]
\end{flushleft}
\subsection*{Page 6}
\begin{flushleft}
\bigskip\textbf{Question 7.} Let f(x) = \sqrtx2 + 1 x.\par
-\par
\textbf{(i)} Using a binomial series, or otherwise, show that, for large x , \sqrtx2 + 1 x + .\par
| | \approx | | 2 x\par
| |\par
\[Sketch the graph y = f(x).\]
\textbf{(ii)} Let g(x) = tan 1 f(x) and, for x = 0, let k(x) = 1 tan 1 .\par
- ( 2 - x\par
\[(a) Show that g(x) + g( x) = 1 \pi.\]
- 2\par
\[(b) Show that k(x) + k( x) = 0.\]
-\par
\[(c) Show that tan k(x) = tan g(x) for x > 0.\]
\[(d) Sketch the graphs y = g(x) and y = k(x) on the same axes.\]
1 0\par
(e) Evaluate k(x) dx and hence write down the value of g(x) dx.\par
0 - 1\par
\bigskip\textbf{Question 8.} (i) Show that\par
1 1 1 1\par
\[zm+1 = z zm + + zm 1 .\]
-\par
- zm+1 - z zm - zm 1\par
( ) ( ) ( - )\par
Hence prove by induction that, for n  1,\par
n\par
1 1 1\par
\[z2n = z z2r 1 + .\]
-\par
\[- z2n - z z2r 1\]
\[( ) r=1 ( - )\]
)\par
1 1\par
\[Find similarly z2n as a product of z + and a sum.\]
\[- z2n z\]
( )\par
\textbf{(ii)} (a) By choosing z = eiθ, show that\par
n\par
\[sin 2nθ = 2 sin θ cos(2r 1)θ .\]
-\par
\[r=1\]
)\par
\[(b) Use this result, with n = 2, to show that cos 2 \pi = cos 1 \pi 1 .\]
5 5 - 2\par
(c) Use this result, with n = 7, to show that cos 2 \pi+cos 4 \pi+cos 8 \pi+cos 16 \pi = 1 .\par
15 15 15 15 2\par
\textbf{(iii)} Show that sin 1 \pi sin 3 \pi + sin 5 \pi = 1 .\par
14 - 14 14 2\par
\end{flushleft}
\subsection*{Page 7}
\begin{flushleft}
\bigskip\textbf{Section B: Mechanics}\par
\bigskip\textbf{Question 9.} In this question, n  2.\par
\textbf{(i)} A solid, of uniform density, is formed by rotating through 360 about the y-axis the\par
\textasciicircum{}\circ\par
\[region bounded by the part of the curve rn - 1y = rn xn with 0 ( x ( r, and the x-\]
-\par
and y-axes.\par
nr\par
Show that the y-coordinate of the centre of mass of this solid is .\par
2(n + 1)\par
\textbf{(ii)} Show that the normal to the curve rn 1y = rn xn at the point rp, r(1 pn) , where\par
-\par
- -\par
\[0 < p ( 1, meets the y-axis at (0, Y ), where Y = r 1 pn  . (\]
- - npn 2\par
) - )\par
\[In the case n = 4, show that the greatest value of Y is 1 r.\]
\textbf{(iii)} A solid is formed by rotating through 360 about the y-axis the region bounded by the\par
\textasciicircum{}\circ\par
\[curves r3y = r^4 x^4 and ry = (r^2 x^2), both for 0 ( x ( r.\]
- - -\par
A and B are the points (0, r) and (0, r), respectively, on the surface of the solid.\par
-\par
Show that the solid can rest in equilibrium on a horizontal surface with the vector -A-\toB\par
at three different, non-zero, angles to the upward vertical. You should not attempt to\par
find these angles.\par
\end{flushleft}
\subsection*{Page 8}
\begin{flushleft}
\bigskip\textbf{Question 10.} A plank AB of length L initially lies horizontally at rest along the x-axis on a flat surface,\par
with A at the origin.\par
\[Point C on the plank is such that AC has length sL, where 0 < s < 1.\]
End A is then raised vertically along the y-axis so that its height above the horizontal surface\par
at time t is h(t), while end B remains in contact with the flat surface and on the x-axis.\par
The function h(t) satisfies the differential equation\par
\[d2h dh\]
\[= ω2h, with h(0) = 0 and = ωL at t = 0,\]
\[dt2 - dt\]
where ω is a positive constant.\par
A particle P of mass m remains in contact with the plank at point C.\par
\textbf{(i)} Show that the x-coordinate of P is sL cos ωt, and find a similar expression for its\par
y-coordinate.\par
\textbf{(ii)} Find expressions for the x- and y-components of the acceleration of the particle.\par
\textbf{(iii)} N and F are the upward normal and frictional components, respectively, of the force\par
of the plank on the particle. Show that\par
\[N = mg(1 k sin ωt) cos ωt,\]
-\par
and that\par
\[F = mgsk + N tan ωt\]
Lω2\par
\[where k = .\]
g\par
\textbf{(iv)} The coefficient of friction between the particle and the plank is tan α, where α is an\par
acute angle.\par
\[Show that the particle will not slip initially, provided sk < tan α.\]
Show further that, in this case, the particle will slip\par
while N is still positive,\par
•\par
when the plank makes an angle less than α to the horizontal.\par
•\par
\end{flushleft}
\subsection*{Page 9}
\begin{flushleft}
\bigskip\textbf{Section C: Probability and Statistics}\par
\bigskip\textbf{Question 11.} (i) Let λ > 0. The independent random variables X , X , . . . , X all have probability\par
1 2 n\par
density function\par
\[f(t) = λe - λt t  0\]
\[0 t < 0\]
and cumulative distribution function F(x).\par
The value of random variable Y is the largest of the values X , X , . . . , X .\par
1 2 n\par
Show that the cumulative distribution function of Y is given, for y  0, by\par
n\par
\[G(y) = 1 e λy .\]
-\par
-\par
( )\par
\textbf{(ii)} The values L(α) and U(α), where 0 < α ( 1 , are such that\par
\[P(Y < L(α)) = α and P(Y > U(α)) = α.\]
Show that\par
\[L(α) = ln 1 α n\]
- λ -\par
( )\par
and write down a similar expression for U(α).\par
\textbf{(iii)} Use the approximation et 1 + t, for t small, to show that, for sufficiently large n,\par
\approx | |\par
λL(α) ln(n) ln ln .\par
\approx - α\par
) ) ((\par
\textbf{(iv)} Hence show that the median of Y tends to infinity as n increases, but that the width\par
of the interval U(α) L(α) tends to a value which is independent of n.\par
-\par
\textbf{(v)} You are given that, for t small, ln(1 + t) t and that e3 20.\par
| | \approx \approx\par
Show that, for sufficiently large n, there is an interval of width approximately 4λ 1 in\par
-\par
which Y lies with probability 0.9 .\par
\end{flushleft}
\subsection*{Page 10}
\begin{flushleft}
\bigskip\textbf{Question 12.} (i) Show that, for any functions f and g, and for any m  0,\par
m+1 m m s+1\par
\[f(r) g(s) = g(s) f(r) .\]
( ) ( )\par
\[r=1 s=r 1 s=0 r=1\]
-\par
\textbf{(ii)} The random variables X , X , X , . . . are defined as follows\par
0 1 2\par
X takes the value 0 with probability 1;\par
•\par
\[X takes the values 0, 1, . . . , X + 1 with equal probability, for n = 0, 1, . . . .\]
n+1 n\par
•\par
(a) Write down E(X ).\par
\[Find P(X = 0) and P(X = 1) and show that P(X = 2) = 1 .\]
2 2 2 6\par
Hence calculate E(X ).\par
(b) For n  1, show that\par
n 1\par
\[- P(X n 1 = s)\]
\[P(X n = 0) = -\]
s + 2\par
\[s=0\]
\[and find a similar expression for P(X = r), for r = 1, 2, . . . , n.\]
n\par
\[(c) Hence show that E(X ) = 1 (1 + E(X )).\]
n 2 n 1\par
-\par
\[Find an expression for E(X ) in terms of n, for n = 1, 2, . . . .\]
n\par
\end{flushleft}
\end{document}
